\section{Related Work}
\label{sec:formatting}

Early work in FER relied on handcrafted features such as Gabor filters, Local Binary Patterns (LBP), and Histogram of Oriented Gradients (HOG), combined with classical classifiers such as Support Vector Machines (SVM). Although these methods offered interpretable feature representations, they struggled with illumination changes, pose variations, and limited generalization to unconstrained real-world scenarios. \cite{Jayaswal, REVINA2021619}

Over the past two decades, facial expression recognition (FER) has undergone a significant paradigm shift, moving from rule-based and handcrafted feature approaches to data-driven, context-aware deep learning models. The introduction of convolutional neural networks (CNNs) marked a major advance in image classification, enabling automatic hierarchical feature learning. Early architectures such as \textit{LeNet} demonstrated the effectiveness of convolutional features in small-scale images \cite{GradientLeNet}, while later models such as \textit{AlexNet} showed that deeper CNNs can achieve significant performance improvements when trained on large datasets \cite{AlexNet}.

Subsequent architectures such as \textit{VGG} and \textit{ResNet} introduced deeper architectures and residual connections, enabling more powerful feature extraction and improved generalization performance \cite{deepconv, deepres}. In the context of FER, previous work suggests that \textit{ResNet}-based architectures can achieve better performance compared to classical CNNs \cite{gao2024cnnresnet}. Other architectures emphasize parameter efficiency. For example, \textit{EfficientNet} introduces a compound scaling strategy that jointly balances network, width, and input resolution \cite{tan2020efficientnet}, while \textit{MobileNet} relies on lightweight convolutional operations to reduce computational complexity and model size \cite{howard2017mobilenet}. Building on these advances, standard and custom CNN architectures have become a dominant methodology for FER \cite{pereira}. They are widely used in the literature, including for low-resolution benchmarks such as the FER2013 dataset \cite{Mollahosseini_2016,khaireddin2021facialemotionrecognitionstate}.

More recently, \textit{Vision Transformers} (ViT) have been proposed as an alternative by modeling global dependencies through self-attention mechanisms. Although ViTs have shown strong performance in various computer vision tasks, they typically require large-scale training data and substantial computational resources, due to a lack of strong inductive biases \cite{dosovitskiy2021vit}. This limitation may be particularly pronounced in FER, where datasets are relatively small. To address these challenges, recent work has explored hybrid CNN–Transformer architectures that combine the local feature extraction capabilities of CNNs with the global context modeling of transformers. While these hybrid models have shown promising results, they introduce additional complexity and raise challenges related to training efficiency and generalization \cite{Haruna_2025, Wu2021convvit}.

